\scanpage{25}
une bijection

\[
\operatorname{Hom}_{\mathrm{Esp}}(\beta I,X)\simeq\operatorname{Hom}_{\mathrm{Préord}}(I,\alpha X)\tag*{$(*)$}
\]

% slanted
\margin{\uline{Cor} le foncteur $\alpha$ $(\mathrm{Esp})\to(\mathrm{Préor})$ commute aux $\varprojlim$ le foncteur $\beta$ aux $\varinjlim$. Mais de plus $\beta$ est \uline{exact à gauche} i.e. sur $X\times Y$ la top. associée au préordre produit est le produit des topologies (Mais $\beta$ ne commute pas aux produits quelconques, un \deleted{foncteur} produit d'espaces topologiques essentiels \unsure{discrets} dont les pts sont essentiels n'est pas néc. essentiel. Ex. produit infini d'espaces top. \uline{discrets}.}

\unsure{$*$ exprime|En effet}: une application

\[
I\xrightarrow{f}X
\]

est croissante sss elle est continue (\unsure{$[$}$X$ étant muni de la
relation d'ordre $\alpha_X(T_X)$, $I$ de la topologie $\beta_I(\omega_I)$).
En effet, si $f$ est continue, on sait qu'elle est croissante pour les
\deleted{topologies}\inserted{ordres} associés aux topologies, or ce sont l'ordre de départ de $I$, et $\uline{\alpha_X(T_X)}$.
Et si $f$ est croissante, \deleted{\unsure{soit ouvert dans $X$}} l'image inverse d'une partie de $X$
ouverte pour $\beta_X\alpha_X(T_X)$ est ouverte dans $I$, mais toute partie
\deleted{\unsure{de $I$}} de $X$ \inserted{ouv\textsuperscript{te}} l'est pour $\beta_X\alpha_{\unsure{\alpha|X}}(T_X)$, donc on a gagné.

e) \uline{Corollaire} Le foncteur

\[
\beta:(\mathrm{Préord})\longrightarrow(\mathrm{Esp})
\]

est pleinement fidèle. Son image essentielle est formée des
espaces $(X,T_X)$ pour lesquels $T_X\in\beta_X(\underline\omega(X))$ ; ce sont aussi
les espaces topologiques pour lesquels tout point $x\in X$ a un
plus petit voisinage ouvert $O_x$. [C'est évidemment nécessaire ;
\deleted{\unsure{et}} pour prouver que c'est suffisant, il faut pr. que

\[
O_x=\{y\in X\mid y\leq_T x\}.
\]

Mais \deleted{\unsure{évident, $O_x$ étant}} \deleted{pour prouver $\subset$, noter que si $y\in O_x$, \unsure{on a} $O_y$}
$x\in\overline{\{y\}}$ signifie évidemt ($O_x$ étant le plus petit \deleted{ouvert}\inserted{voisinage} de $x$)
que $O_x\ni y$, OK.

\margin{NB les topologies obtenues sont aussi celles pour lesquelles toute intersection d'ouverts est ouverte [NB Ce n'est pas stable par passage à l'espace sobre associé ....]}

f) Soit \deleted{\unsure{un}} $X,T$, alors $X$ est présobre $(\overline{\{x\}}=\overline{\{y\}}\ \Rightarrow\ x=y)$
\deleted{\unsure{g)}} sss \deleted{\unsure{le}} préordre $\omega(T)$ est un \uline{ordre}. \deleted{\unsure{$X$}},

\deleted{g) Soit $X$ muni d'un préordre, d'où topologie associée. Désignons par $\underline X$ la catégorie associée à $X$ ordonné ($\operatorname{Hom}(x,y)=\begin{cases}\varnothing&\text{si }\deleted{\unsure{y}}\,x\nleq y\\\{e\}&\text{si }\deleted{\unsure{x}}\,x\leq y\end{cases}$ i.e. $\operatorname{Hom}(x,y)\ne\varnothing$ sss $O_y\supset O_x$) Montrer que le \deleted{foncteur}}

\deleted{$F\longmapsto(x\mapsto F(O_x))$ de $\operatorname{Top}(X)$ dans $\widehat{\underline X}$ est une équivalence de cat.}

\[
\deleted{\operatorname{Top}(X)\simeq\widehat{\underline X}}
\]

\deleted{On identifie les $x\in X$ à des sous-objets de $e_{\widehat{\underline X}}$, et les cribles de $X$ (i.e. \inserted{des ouverts de $\widetilde X$} ouverts de $X$) à des s-objets de $e_{\widehat{\underline X}}$, i.e.}
% archive sheet number: 24
